\documentclass[letterpaper, onecolumn, draftclsnofoot, 10pt, compsoc]{IEEEtran}
\usepackage{graphicx}
\usepackage{url}
\usepackage{setspace}

\usepackage{geometry}
\geometry{textheight=9.5in, textwidth=7in}

% 1. Fill in these details
%\def \CapstoneTeamName{		The Cleverly Named Team}
\def \CapstoneTeamNumber{		63}
\def \GroupMemberOne{			Corynna Park}
\def \GroupMemberTwo{			Joshua Lioy}
\def \GroupMemberThree{			Jackson Wells}
\def \CapstoneProjectName{		Molecules in 3D?! And in color!? That I can hold in my hand? No way!!!}
\def \CapstoneSponsorCompany{	Oregon State University College of Science Department of Biochemistry and Biophysics}
\def \CapstoneSponsorPerson{		Dr. Victor Hsu}

% 2. Uncomment the appropriate line below so that the document type works
\def \DocType{		Problem Statement
				%Requirements Document
				%Technology Review
				%Design Document
				%Progress Report
				}
			
\newcommand{\NameSigPair}[1]{\par
\makebox[2.75in][r]{#1} \hfil 	\makebox[3.25in]{\makebox[2.25in]{\hrulefill} \hfill		\makebox[.75in]{\hrulefill}}
\par\vspace{-12pt} \textit{\tiny\noindent
\makebox[2.75in]{} \hfil		\makebox[3.25in]{\makebox[2.25in][r]{Signature} \hfill	\makebox[.75in][r]{Date}}}}
% 3. If the document is not to be signed, uncomment the RENEWcommand below
%\renewcommand{\NameSigPair}[1]{#1}

%%%%%%%%%%%%%%%%%%%%%%%%%%%%%%%%%%%%%%%
\begin{document}
\begin{titlepage}
    \pagenumbering{gobble}
    \begin{singlespace}
    	\includegraphics[height=4cm]{coe_v_spot1}
        \hfill 
        % 4. If you have a logo, use this includegraphics command to put it on the coversheet.
        %\includegraphics[height=4cm]{CompanyLogo}   
        \par\vspace{.2in}
        \centering
        \scshape{
            \huge CS Capstone \DocType \par
            {\large\today}\par
            \vspace{.5in}
            \textbf{\Huge\CapstoneProjectName}\par
            \vfill
            {\large Prepared for}\par
            \Huge \CapstoneSponsorCompany\par
            \vspace{5pt}
            {\Large\NameSigPair{\CapstoneSponsorPerson}\par}
            {\large Prepared by }\par
            Group\CapstoneTeamNumber\par
            % 5. comment out the line below this one if you do not wish to name your team
            %\CapstoneTeamName\par 
            \vspace{5pt}
            {\Large
                %\NameSigPair{\GroupMemberOne}\par
                \NameSigPair{\GroupMemberTwo}\par
                %\NameSigPair{\GroupMemberThree}\par
            }
            \vspace{20pt}
        }
        \begin{abstract}
        % 6. Fill in your abstract    
        	The overall goal of our project is to design a new workflow system for printing multicolored 3d representations of molecules, specifically biomolecules such as proteins and DNA structures. The existing workflow system for creating models of biomolecules in the OSU Department of Biochemistry and Biophysics consists of printing of a white object and painting it by hand. This is of course a slow and inefficient use time and teaching resources.\par 
            This problem may be addressed by implementing a 3d printing system that makes use of multiple colored filaments, in our case up to four different colors, that creates a multi-colored object. Our solution will need to make use of the Mosaic Manufacturing Palette+, open source or freely available software tools, and custom algorithms and or programs to create an appropriately tailored and streamlined workflow.\par
            At the end of the project we plan to have a workflow system that will allow a worker who has a basic knowledge of 3d printing to create a multi-colored object just as easily as a monochrome object.
        \end{abstract}     
    \end{singlespace}
\end{titlepage}
\newpage
\pagenumbering{arabic}
\tableofcontents
% 7. uncomment this (if applicable). Consider adding a page break.
%\listoffigures
%\listoftables
\clearpage

% 8. now you write!
\section{Project Definition}
    \begin{singlespace}
        The Oregon State University College of Science Department of Biochemistry and Biophysics (DBB) currently lacks the capability to create multi-colored 3d objects with their existing 3d printing workflow. The DBB makes use of 3d printed representations of biomolecular structures such as proteins and DNA structures so that students can have a visual aid during their learning experience. Presently, for the DBB to create a 3d representation of a biomolecular structure, one must print off a monochromatic object and then proceed to paint it by hand so that they are able to use different colors to represent different molecules. This creates an additional problem in that there are a limited number of objects that can be created for a given class. This limits the effectiveness of the 3d printed biomolecular representations as teaching tools.\par
        The other component that needs to be addressed is the hardware and software tools that are currently being used by the DBB. The current 3d printing hardware that is available to the DBB only allows for the use of one colored filament when printing an object. As for the software, the DBB can create 3d printing data files that show multi-colored objects on the computer, but they do not have a way to transfer this color data to a 3d printer. Thus, there us currently no workflow in place for the creation of 3d printed biomolecular representations for educational purposes.
    \end{singlespace}

\section{Proposed Project Solution}
    \begin{singlespace}
        To address the problem at the hardware level, the DBB has already purchased a multi-filament attachment from Mosaic Manufacturing called the Palette+. The Palette+ will allow us to make use of up to four distinct colored filaments to print the 3d biomolecular representations. The Palette+ will attach to an existing 3d printer that extrudes only a single filament. The Palette+ will handle the selection and feeding of the various colored filaments while the 3d printer will perform in the same manner as it did before the addition of the Palette+.\par 
        To address the problem at a software level there are several tools that could potentially be used to implement the necessary workflow. Firstly, making use of existing open source or freely available tools to input usable color data into the 3d object files. These files then need to be “sliced” so that the object data can be interpreted by the 3d printer. This is done by converting the object files into gcode, which is a language that the 3d printer uses to print a given object.\par
        At this point in the process, the gcode needs to be processed so that the Palette+ can determine which filament to feed to the 3d printer at any given time. This is where it may be appropriate for us to create a proprietary tool that interprets the gcode and sends the correct commands to the Palette+. This could be done through scripting or other means depending on what fits best in to the workflow that we design and implement.\par
    \end{singlespace}
    
\section{Performance Metrics}
    \begin{singlespace}
        The first goal of the team should be to familiarize ourselves with the 3d printing technology that we will be utilizing throughout the course of the project. This includes the two 3d printers in the DBB lab as well as the Palette+ module that will be the crux of the project. We will need to demonstrate to Dr. Victor Hsu that we have an appropriate grasp of the basic 3d printing technology (i.e. our ability to print monochromatic objects), this will include familiarization of the existing workflow and software that is being used in the DBB.\par
        The second goal of the team should be the implementation of a streamlined workflow. This will need to entail the use of both the existing hardware resources as well as a preliminary selection of software tools, open source or ones that we create, that will be used in the workflow.\par
        The third goal of the team should be the implementation of our workflow design once it has been approved by Dr. Victor Hsu. The implementation will need to take place in two phases. Phase one will consist of putting in place the basic structure of the workflow so that we have a fully functional base that we will be able to improve upon later. Phase two will consist of refining and streamlining the workflow so that a person with limited knowledge of 3d printing can make use of the equipment in the DBB to create a multi-colored 3d object.\par 
    \end{singlespace}
\end{document}